\clearpage
\section{SneakySnake prefiltering}
\label{sec:prefilter-role}
\label{sec:prefilter-sneakysnake}

A plausible chain can still contain too many edits to satisfy the query.
A prefilter tries to reject such candidates before spending time on alignment.
BioseqDB incorporates the existing SneakySnake algorithm as a component of
this step~\cite{sneakysnake_paper}. The algorithm described here is borrowed
from that work, rather than a new filtering method developed for this thesis.

For an edit limit \(E\), SneakySnake compares the sequences at shifts from
\(-E\) to \(E\). At the current position it follows the longest matching
run across these shifts, crosses one mismatch, and repeats. It counts these
crossings but does not charge for switching between shifts. Its estimate can
therefore be lower than the true edit distance. If the estimate exceeds
\(E\), the candidate can be rejected. Otherwise it still needs alignment.
The filter can save time because it compares runs of equal symbols without
computing a full matrix of alignment scores or reconstructing a path.

BioseqDB applies the test to mandatory intervals between anchors. It leaves
the optional outer extensions to the aligner, which may stop before their
ends. For unequal-length intervals, the current implementation uses only the
lower bound \(\bigl||Q|-|T|\bigr|\). Surviving candidates are aligned and
checked against the original edit predicate, such as
\(\texttt{bioseq.edit(seq)}<10\). Passing the filter alone does not establish
that this condition holds.

% Existing filter illustrations retained as source material.
% \begin{figure}[H]
% \centering
% \begingroup
% \small
% \setlength{\tabcolsep}{4pt}
% \renewcommand{\arraystretch}{1.2}
% \newcommand{\mazeMatch}{\cellcolor{white}\(\cdot\)}
% \newcommand{\mazeObstacle}{\cellcolor{black!70}\textcolor{white}{\(\times\)}}
% \newcommand{\mazeRoute}{\cellcolor{orange!35}\textcolor{orange!80!black}{\(\rightarrow\)}}
% \newcommand{\mazeCross}{\cellcolor{orange!70}\(\times\)}
% \begin{tabular}{c|*{12}{c}}
%  & 0 & 1 & 2 & 3 & 4 & 5 & 6 & 7 & 8 & 9 & 10 & 11\\
% \hline
% \(d=+2\) & \mazeMatch & \mazeMatch & \mazeMatch & \mazeObstacle &
% \mazeMatch & \mazeMatch & \mazeMatch & \mazeObstacle & \mazeMatch &
% \mazeMatch & \mazeObstacle & \mazeObstacle\\
% \(d=+1\) & \mazeMatch & \mazeMatch & \mazeMatch & \mazeObstacle &
% \mazeMatch & \mazeMatch & \mazeMatch & \mazeObstacle & \mazeMatch &
% \mazeMatch & \mazeMatch & \mazeObstacle\\
% \(d=0\) & \mazeRoute & \mazeRoute & \mazeRoute & \mazeCross &
% \mazeRoute & \mazeRoute & \mazeRoute & \mazeCross & \mazeRoute &
% \mazeRoute & \mazeRoute & \mazeCross\\
% \(d=-1\) & \mazeObstacle & \mazeMatch & \mazeMatch & \mazeObstacle &
% \mazeMatch & \mazeMatch & \mazeMatch & \mazeObstacle & \mazeMatch &
% \mazeMatch & \mazeMatch & \mazeObstacle\\
% \(d=-2\) & \mazeObstacle & \mazeObstacle & \mazeMatch & \mazeObstacle &
% \mazeMatch & \mazeMatch & \mazeMatch & \mazeObstacle & \mazeMatch &
% \mazeMatch & \mazeMatch & \mazeObstacle\\
% \end{tabular}

% \vspace{0.4em}
% \begin{tabular}{clclcl}
% \cellcolor{white}\(\cdot\) & match &
% \cellcolor{black!70}\textcolor{white}{\(\times\)} & obstacle &
% \cellcolor{orange!35}\(\rightarrow\) & selected escape segment
% \end{tabular}
% \endgroup
% \caption{A rejecting SneakySnake chip-maze walk for
% \(Q=\mathtt{AAAAAAAAAAAA}\), \(T=\mathtt{AAACAAACAAAC}\), and \(E=2\).
% Each row represents a query-to-target displacement. Amber cells show the
% successive escape segments and three crossed obstacles on the unshifted row.}
% \label{fig:prefilter-chip-maze}
% \end{figure}

% \begin{figure}[H]
% \centering
% \small
% \renewcommand{\arraystretch}{1.15}
% \begin{tabular}{c}
% \fbox{\parbox{0.74\textwidth}{\centering
% Construct the unshifted mismatch mask \(H_0\)}}\\[0.35em]
% \(\operatorname{popcount}(H_0)\leq E\): retain \qquad
% \(E=0\): reject by exact Hamming comparison\\[0.35em]
% \(\Downarrow\) otherwise\\[0.35em]
% \fbox{\parbox{0.74\textwidth}{\centering
% Construct the \(2E\) shifted masks, amend each mask independently,\\
% then set the boundary-adjacent compensation bits}}\\[0.35em]
% \(\Downarrow\)\\[0.35em]
% \fbox{\parbox{0.74\textwidth}{\centering
% Intersect the masks and compute the capped four-bit group estimate}}\\[0.35em]
% \(\Downarrow\) retain when the estimate is at most \(E\)\\[0.35em]
% \fbox{\parbox{0.74\textwidth}{\centering
% For a proposed rejection, evaluate exact thresholded edit distance}}\\[0.35em]
% \(D(Q,T)\leq E\): retain \qquad \(D(Q,T)>E\): reject
% \end{tabular}
% \caption{GateKeeper decision flow in BioSeqDB. The unshifted mask provides an
% exact rejection certificate when \(E=0\). For positive thresholds, shifted
% masks can retain a candidate or propose rejection, and the exact threshold
% guard confirms the proposal.}
% \label{fig:prefilter-gatekeeper}
% \end{figure}
